%The next statement creates the title page.
\frame{\titlepage}

%This block of code is for the table of contents after
%the title page
\begin{frame}
\frametitle{Table of Contents}
\tableofcontents
\end{frame}


\section{Détection de $f_{0}$ pour deux signaux \texttt{fluteircam.wav} et \texttt{voiceP.wav}}
\subsection{Détection de $f_{0}$ avec la FFT}

\begin{frame}
\frametitle{Sample frame title}
This is a text in second frame. For the sake of showing an example.

\begin{itemize}
    \item<1-> Text visible on slide 1
    \item<2-> Text visible on slide 2
    \item<3> Text visible on slides 3
    \item<4-> Text visible on slide 4
\end{itemize}
\end{frame}

\subsection{Détection de $f_{0}$ avec l'autocorrélation}

\begin{frame}
    
\end{frame}

\subsection{Comparaison des performances}

\begin{frame}
    
\end{frame}

\subsection{Détection du silence}

\begin{frame}
    La détection de l'énergie se fait à l"aide de l'énergie.
\end{frame}

\subsection{Détection du tempo}

\begin{frame}
    
\end{frame}

\subsection{Détection du vibrato}

\begin{frame}
    
\end{frame}

\section{Croisement de deux sinus}

\section{Description de \texttt{myson.wav}}